\section{Canonical gain surfaces and defect flags}
\label{sec:canonical-gain-surface}

M\"uller and Taktikos propose separating the quality of certain $abc$
equations into an approximation gain and a power gain
\cite{MuellerTaktikos2026}.  The factorization itself is exact and useful, but
the proposed fixed bounds on the two factors do not supply the coefficient-one quantifier
in the standard $abc$ conjecture.  We record the precise correlation that is
needed and then introduce a finite-dimensional defect flag on which different
arithmetic mechanisms can act.

The source requires two cautions.  Definition~2.6 prints
$\log(q_n^s k)$ in the numerator when $d_n>0$, whereas the proof of
Theorem~2.7 uses $\log(p_n^3)$ for the same sign.  These are distinct when
$p_n^3=kq_n^3+d_n$ and $d_n>0$.  Moreover, the printed proof of
Theorem~2.8 obtains a limiting equality from a one-sided asymptotic estimate
and refers to a finite verification below an unspecified cutoff.  We therefore
use no continued-fraction assertion from that theorem.  The frozen source and
its audit extraction are recorded under
\path{research/sources/abc_gain_surface_2026_09_03/}.

\subsection{The gain surface}

Let $H,N,\mathcal R>1$ and write
\[
 h=\log H,\qquad n=\log N,\qquad \rho=\log \mathcal R.
\]
Define
\[
 q(H,\mathcal R)=\frac{h}{\rho},\qquad
 A(H,N)=\frac{h}{n},\qquad
 P(N,\mathcal R)=\frac{n}{\rho}.
\]

\begin{proposition}[Exact gain surface]
\label{prop:exact-gain-surface}
For every intermediate scale $N>1$,
\[
q(H,\mathcal R)=A(H,N)P(N,\mathcal R).
\]
\end{proposition}

\begin{proof}
Both $n$ and $\rho$ are positive, and
$(h/n)(n/\rho)=h/\rho$.
\end{proof}

The level set $q=1+\varepsilon$ is the hyperbola
$AP=1+\varepsilon$.  Separate bounds $A\le A_0$ and $P\le P_0$
give only $q\le A_0P_0$.  Thus the proposed constants $3/2$ and $3$
give a weak exponent $9/2$, rather than the standard $1+\varepsilon$
for every positive $\varepsilon$.

This failure is not an issue of endpoints.  In logarithmic coordinates take
$(\rho,h,n)=(2,4,5)$ and exponentiate.  Then $1<\mathcal R<H<N$ while
\[
 A=\frac45<\frac32,\qquad P=\frac52<3,
 \qquad q=2>\frac32.
\]
This is a full-premise countermodel to the abstract implication from the two
independent bounds.  It is not an $abc$ triple.

\subsection{The canonical arithmetic corridor}

Let $a+b=c$ be a primitive positive triple with $a,b>1$, and set
\[
 M=abc,\qquad \mathcal R=\rad(abc),\qquad
 A_{\rm can}=\frac{\log c}{\log M},\qquad
 P_{\rm can}=\frac{\log M}{\log \mathcal R}.
\]

\begin{theorem}[Canonical gain corridor]
\label{thm:canonical-gain-corridor}
Every such triple satisfies
\[
 c^2<M<c^3,
 \qquad
 \frac13<A_{\rm can}<\frac12,
 \qquad
 q(a,b,c)=A_{\rm can}P_{\rm can}.
\]
Consequently
\[
 \frac13P_{\rm can}<q(a,b,c)<\frac12P_{\rm can}.
\]
\end{theorem}

\begin{proof}
Because $a,b>1$ and $(a,b)=1$, the pair cannot be $(2,2)$, so
$(a-1)(b-1)>1$.  Expansion and $a+b=c$ give $ab>c$, whence
$abc>c^2$.  Also $a<c$ and $b<c$, so $abc<c^3$.  Strict monotonicity
of the logarithm gives
\[
 2\log c<\log M<3\log c,
\]
which is equivalent to the claimed corridor.  The factorization is
Proposition~\ref{prop:exact-gain-surface}; multiplying the corridor by the
positive number $P_{\rm can}$ proves the last display.
\end{proof}

\begin{corollary}[Necessary canonical power concentration]
\label{cor:canonical-power-concentration}
If $q(a,b,c)\ge1+\varepsilon$, then
\[
 P_{\rm can}>2(1+\varepsilon).
\]
Conversely, $P_{\rm can}\le2(1+\varepsilon)$ implies
$q(a,b,c)<1+\varepsilon$.
\end{corollary}

\begin{proof}
Use $q<P_{\rm can}/2$ from
Theorem~\ref{thm:canonical-gain-corridor}.
\end{proof}

The coefficient two is the same sharp threshold found by the symmetric
product transfer, but the factorized form exposes the cancellation that the
coefficient statement hides.

\subsection{The affine defect plane}

Put $h=\log c$, $m=\log M$, and $\rho=\log \mathcal R$, and define
\[
 X=\frac{m-\rho}{\rho},\qquad
 Y=\frac{m-h}{\rho}.
\]
Here $X$ is the normalized multiplicity discarded by the radical and $Y$ is
the normalized product-to-height slack.

\begin{theorem}[Exact cancellation]
\label{thm:gain-defect-cancellation}
One has
\[
 \boxed{q(a,b,c)=1+X-Y.}
\]
Consequently the effective logarithmic inequality
\[
 h\le(1+\varepsilon)\rho+C
\]
is equivalent to
\[
 X-Y\le\varepsilon+\frac{C}{\rho}.
\]
\end{theorem}

\begin{proof}
Cancellation gives
\[
 1+\frac{m-\rho}{\rho}-\frac{m-h}{\rho}=\frac{h}{\rho}.
\]
The second assertion follows by subtracting one and dividing by $\rho>0$.
\end{proof}

Thus a large power excess can be harmless when the approximation slack is
comparably large.  Independent constant bounds erase this relation.

\subsection{A scale groupoid and higher-dimensional flags}

On real scale coordinates define
\[
 \delta(x,y)=y-x.
\]
This is an additive one-cocycle on the pair groupoid:
\[
 \delta(x,x)=0,
 \qquad
 \delta(x,z)=\delta(x,y)+\delta(y,z).
\]
On nonzero coordinates, $g(x,y)=y/x$ is the corresponding multiplicative
cocycle.  Both identities are immediate by cancellation.

A \emph{defect flag} is a finite path
\[
 \rho=s_0,s_1,\ldots,s_k=m,
 \qquad d_i=s_i-s_{i-1}.
\]
Its total is path-independent:
\[
 \sum_{i=1}^k d_i=m-\rho.
\]
In particular every closed scale flag has trivial additive and multiplicative
holonomy.  Refining a factorization cannot manufacture a saving; arithmetic
must restrict the admissible paths or bound their costs.

\begin{theorem}[Defect-budget half-space]
\label{thm:defect-budget-half-space}
Suppose
\[
 \sum_i d_i\le\sum_i B_i
\]
and
\[
 \sum_i B_i\le(m-h)+\varepsilon \rho+C.
\]
Then
\[
 h\le(1+\varepsilon)\rho+C.
\]
\end{theorem}

\begin{proof}
The cocycle law telescopes, so
$m-\rho=\sum_i d_i\le\sum_iB_i$.  Combining this with the second budget and
cancelling $m$ gives the result.
\end{proof}

For a fixed total allowance and no sign constraints, the vectors
$(B_1,\ldots,B_k)$ lie in an affine half-space.  Its intersection with the
nonnegative orthant is a simplex when the allowance is nonnegative.  Packet
compression, Pell transversality, IUT log-volume estimates, and analytic
radical bounds may therefore be assigned to different coordinates, but their
total must remain below the approximation slack plus the
$\varepsilon$-allowance.

\subsection{An exact counterexample boundary}

The tempting universal strengthening $P_{\rm can}\le3$ is false.  For the
primitive hit
\[
 3+125=128
\]
one has
\[
 M=48000=2^7\cdot3\cdot5^3,
 \qquad \mathcal R=30,
\]
and $48000>30^3=27000$.  Hence
\[
 P_{\rm can}=\frac{\log48000}{\log30}>3.
\]
This retires only the universal power-three cap.  It does not refute the
correlated defect barrier or the $abc$ conjecture.

An exhaustive computation of the $5{,}465{,}583$ unordered primitive
nonunit triples with $c\le6000$ found no violation of the canonical
approximation corridor, $65$ triples of quality above one, and $14$ triples
with canonical power gain above three.  These data are finite evidence only.

The exact open gate is, for every $\varepsilon>0$, to prove
\[
 \sum_i d_i\le(m-h)+\varepsilon \rho+C_\varepsilon.
\]
This is equivalent to standard $abc$, so the coordinate construction alone
does not close the conjecture.  It supplies a common, noncircular accounting
surface on which the different routes can be combined.
